\chapter{The obstacle map}\label{chap:obstacles}

The preceding chapters closed the repeated-root sector (Theorem~$7'$, Chapter~\ref{chap:repair}) and located the remaining open content of Theorem~7 at a single open statement, Definition~\ref{Conjecture1_normal}.  This chapter records
what stands between the present state of knowledge and that statement.  Each wall
below is documented by a machine-verifiable artefact (a kernel check, a two-system
computer-algebra record with checksum, or a live computation log), together with a
\emph{reopening gate}: the explicit external event that would dissolve it.  The
purpose is the one served by the dependency graph throughout this blueprint: to
draw the boundary between the settled and the open so that it can be audited, and
so that future work --- ours or anyone's --- knows exactly where to push.
None of the walls is an assertion of impossibility.

\section{Formalisation ceiling}

\begin{remark}[Wall W1: the proof-assistant library]\label{wall-mathlib}
The Mathlib library, at the pinned commit used throughout this development,
contains no formalisation of Mordell--Weil rank, of Selmer groups, of Chabauty's
method, or of Coleman integration (no declarations for these notions were found by
search of the library source at the pin).  Consequently any completion of
Definition~\ref{Conjecture1_normal} that proceeds through the arithmetic of
elliptic surfaces --- the only known route, by Chapter~\ref{chap:boundary} --- cannot
at present be carried to a kernel-checked proof without disclosed axioms
encapsulating the computer-algebra rank computations.  This is a property of the
current library, not of the mathematics.
\emph{Status (2026-08):} for the repeated-root sector this wall was \emph{bypassed} rather than breached: Theorem~$7'$ (Chapter~\ref{chap:repair}) is rank-free --- both directions are polynomial identities --- so the classification of that sector consumed no Mordell--Weil input and closed with the standard three axioms. The wall remains standing exactly for Definition~\ref{Conjecture1_normal}.
\emph{Reopening gate:} Mathlib acquires Mordell--Weil rank and descent machinery.
\end{remark}

\section{Arithmetic walls}

\begin{remark}[Wall W2: five measured fibres are the noise floor]\label{wall-beta}
In the companion programme attacking Definition~\ref{Conjecture1_normal} through
the fibred structure of the associated elliptic surfaces, the target parity datum
$\beta$ has been measured for exactly five fibres ($n = 2,\dots,6$).  A systematic
mining run over $488$ atomic arithmetic features (with a pre-declared overfit
guard) found $16$ exact matches against $15.25$ expected by chance: agreement over
five fibres is statistically vacuous, and every surviving rule was verified to be
an avatar of the single invariant $[3 \nmid n]$.  Candidate softer structures were
refuted by extension testing.  The mining record, with checksums, is preserved in the companion
programme's archive, outside this repository (cf.~Wall~W5).  \emph{Reopening gate:} measurement of $\beta$ for any fibre with
$n \ge 7$.
\end{remark}

\begin{remark}[Wall W3: the cost of one more fibre]\label{wall-descent}
The reason $n \ge 7$ is unmeasured is concrete.  The $2$-descent on the genus-$2$
Jacobian attached to the fibre $n = 6$ --- an integral model with leading
coefficient of $171$ decimal digits --- has, at the time of writing, been running
continuously for more than three and a half days on a dedicated machine inside the
class-group stage, under an unconditional (non-GRH) regime.  The launch record and
log are preserved.  This is the physical form of Wall~W2: the per-fibre cost of
unconditional descent grows past the practical horizon exactly where finer
structure would first become visible.
\emph{Reopening gate:} completion of the $n = 6$ computation; or an improved model
reducing the coefficient blow-up; or an algorithmic advance in unconditional
class-group computation.
\end{remark}

\begin{remark}[Wall W4: no geometric mechanism, no descent finiteness]\label{wall-covers}
Chapter~\ref{chap:boundary} shows that away from the (rationally empty) symmetric
locus, the conclusion of Theorem~7 demands the lifting of rational points through
one of three nontrivial double covers governed by the single polynomial $Q_4$, and
that the relevant twist classes on the surface admit no finiteness that a descent
argument could exploit.  Every mechanism visible at the function-field level is
therefore closed, and the field-dependence exhibited by the two refutations
(Theorem~\ref{not_T9_bridge_Q} over $\mathbb{Q}$,
Theorem~\ref{not_T0_claim_K} over $\mathbb{Q}(\sqrt{3441})$, both funnelling
through the square class $209$) constrains any future proof to depend essentially
on the arithmetic of the base field.
\emph{Reopening gate:} a new arithmetic mechanism --- for instance a Brauer--Manin
computation on the surface confining its rational points --- or
Definition~\ref{Conjecture1_normal} itself.
\end{remark}

\section{The measured distance}

\begin{remark}[Wall W5: the disclosed-axiom footprint]\label{wall-footprint}
The companion programme --- whose artifacts live outside this repository; the present artifact is std-3 throughout --- has a current summit theorem for
Definition~\ref{Conjecture1_normal} that carries, by independent kernel measurement, an
axiom footprint consisting of the three standard axioms of Lean's classical logic
together with five disclosed axioms: three encapsulating specific computer-algebra
rank computations, one packaging the tail of an even-orbit Mordell--Weil argument,
and one a coordinate function pending replacement by the library's Weierstrass
group law.  This footprint is the honest, machine-readable measure of the distance
between the present state and a complete proof; two of the five are engineering
debts with identified repairs, three are permanent under Wall~W1.
\emph{Reopening gate:} discharge of the load-bearing reduction through
Definition~\ref{Conjecture1_normal} (Chapter~\ref{chap:boundary}); the two
engineering repairs; and, for the remaining three, Wall~W1.
\end{remark}

\section*{Reading the map}

Twenty-six years ago the material of Chapter~\ref{chap:transcript} was printed as a
theorem, and the boundary drawn in this blueprint was invisible.  The walls above
are offered in the opposite spirit: as a verifiable record of where the open
problem actually stands, in a form that a reader --- or a machine --- can audit,
attack, and, we hope, eventually demolish.
